\chapter{The area formula from metric pieces}\label{sec:euclidean-area-metric-proof}
We expand the three cases in \cite{Sim18}, Chapter 2, Theorem 3.3 and its proof, printed pp.~57--60. The proof uses differentiability, countably many pieces on which distances are close to a fixed linear metric, and an injective graph regularization of the critical set. It uses the normalization $\mathcal H^m=\mathcal L^m$ on $\mathbb R^m$, the diameter-cover Lipschitz estimate and ordinary linear change of variables. It does not use rectifiable-set tangent theorems, Whitney extension, or coarea. Throughout $1\leq m\leq n$.

\begin{lemma}[An injective linear map determines a positive domain metric]\label{lem:area-positive-metric-factor}
Let $L:\mathbb R^m\to\mathbb R^n$ be injective and linear, and put $T=(L^*L)^{1/2}$. Then $T$ is positive definite, $|Tv|=|Lv|$ for every $v$, and $\det T=J_mL$. For every Lebesgue-measurable $E\subset\mathbb R^m$,
\[
 \mathcal H^m(L(E))=\mathcal H^m(T(E))=J_mL\,\mathcal L^m(E).
\]
This is the linear normalization used in \cite{Sim18}, proof of Theorem 3.3, pp.~58--59.
\end{lemma}
\begin{proof}
The spectral theorem constructs the positive square root and gives $|Tv|^2=\langle L^*Lv,v\rangle=|Lv|^2$. Its determinant is the square root of $\det(L^*L)$. The map $LT^{-1}$ is an isometric linear embedding, so the diameter-cover definition gives invariance of Hausdorff measure under it. On the domain plane, the normalization recorded in Definition \ref{def:hausdorff-measure} identifies $\mathcal H^m$ with Lebesgue measure; ordinary linear change of variables gives $\mathcal L^m(T(E))=|\det T|\mathcal L^m(E)$. This establishes the linear case without an area formula for nonlinear maps. Infinite values follow by increasing bounded pieces.
\end{proof}

\begin{lemma}[Measurability of a finite-volume Lipschitz image]\label{lem:finite-volume-lipschitz-image-measurable}
Let $f:\mathbb R^m\to\mathbb R^n$ be Lipschitz, and let $E\subset\mathbb R^m$ be Lebesgue measurable with finite measure. Then $f(E)$ is $\mathcal H^m$-measurable. If $|E|=0$, then $\mathcal H^m(f(E))=0$.
\end{lemma}
\begin{proof}
Choose compact sets $K_j\subset E$ with $|E\setminus\bigcup_jK_j|=0$ by inner regularity. Images of the compact sets are compact, so their union is Borel. Proposition \ref{prop:lipschitz-measure} gives
\[
 \mathcal H^m(f(Z))\leq (\operatorname{Lip}f)^m\mathcal H^m(Z)=0
 \quad\text{when }|Z|=0.
\]
Thus $f(E)$ differs from the Borel union $\bigcup_jf(K_j)$ by a Hausdorff-null set and is measurable for the complete Hausdorff measure. This also proves the null assertion. The argument extends to arbitrary measurable $E$ by a countable bounded exhaustion. It requires neither injectivity nor any assertion that every continuous image of a Borel set is Borel.
\end{proof}

\begin{proposition}[A countable partition with controlled linear metric]\label{prop:area-full-rank-metric-pieces}
Let $f:\mathbb R^m\to\mathbb R^n$ be Lipschitz, let $A$ be a measurable subset of its differentiability set on which $Df$ has rank $m$, and fix $0<\varepsilon<1$. There is a countable disjoint measurable partition $A=\bigcup_\ell E_\ell$ and positive definite symmetric maps $T_\ell:\mathbb R^m\to\mathbb R^m$ such that, for every $x,x'\in E_\ell$,
\[
 (1-\varepsilon)|T_\ell(x-x')|
 \leq |f(x)-f(x')|
 \leq(1+\varepsilon)|T_\ell(x-x')|,
\]
and at every $x\in E_\ell$,
\[
 (1-\varepsilon)\det T_\ell\leq J_mDf(x)
 \leq(1+\varepsilon)\det T_\ell.
\]
In particular $f$ is injective on each piece. This is the quantified partition underlying \cite{Sim18}, pp.~58--59.
\end{proposition}
\begin{proof}
Choose a countable dense family of injective linear maps $L_j:\mathbb R^m\to\mathbb R^n$ and put $T_j=(L_j^*L_j)^{1/2}$. At a fixed full-rank differentiability point $x$, the smallest singular value of $Df(x)$ is positive. Differentiability makes $|f(x+v)-f(x)-Df(x)v|/|v|$ arbitrarily small for all sufficiently small nonzero $v$. Choose $L_j$ close enough to $Df(x)$ in operator norm and its Jacobian close enough to $J_mDf(x)$. These two facts give the displayed distance and Jacobian inequalities with $T_j$, for all $0<|v|<1/i$ for some positive integer $i$. The positive least singular value converts the small error measured by $|v|$ into an error relative to $|L_jv|=|T_jv|$.

For fixed $j,i$, the set of qualifying $x$ is measurable. The Jacobian condition uses a measurable derivative representative; the distance inequalities can be tested only on rational nonzero vectors $v$ of length less than $1/i$. Continuity in $v$ then supplies the inequalities for all such vectors. This is a countable family of inequalities between measurable functions of $x$. The union of these qualifying sets covers $A$. Assign a point to the first pair $(j,i)$ in a fixed enumeration, and then partition each assigned set by a countable disjoint grid of sets of diameter strictly less than $1/i$. Both distance inequalities now apply to every pair in the resulting piece. Use $T_\ell=T_j$ for that piece. Empty pieces may be omitted. The lower bound and invertibility of $T_j$ imply injectivity. No area estimate has been used to construct this partition.
\end{proof}

\begin{lemma}[Area bounds on one metric piece]\label{lem:area-one-metric-piece}
Suppose $E$ has finite Lebesgue measure and the two bounds in Proposition \ref{prop:area-full-rank-metric-pieces} hold on $E$ with one positive definite $T$. Then
\[
 \frac{(1-\varepsilon)^m}{1+\varepsilon}\int_EJ_mDf\,dx
 \leq\mathcal H^m(f(E))
 \leq\frac{(1+\varepsilon)^m}{1-\varepsilon}\int_EJ_mDf\,dx.
\]
Both comparison factors tend to one as $\varepsilon\downarrow0$.
\end{lemma}
\begin{proof}
The map $h=f\circ T^{-1}:T(E)\to f(E)$ is injective and has Lipschitz constant at most $1+\varepsilon$; its inverse has Lipschitz constant at most $(1-\varepsilon)^{-1}$. Proposition \ref{prop:lipschitz-measure}, applied separately to these two maps on their domains, gives
\[
 (1-\varepsilon)^m\mathcal H^m(T(E))
 \leq\mathcal H^m(f(E))
 \leq(1+\varepsilon)^m\mathcal H^m(T(E)).
\]
Lemma \ref{lem:area-positive-metric-factor} identifies the domain measure as $\det T\,|E|$. Integrate the pointwise bounds $(1-\varepsilon)\det T\leq J_mDf\leq(1+\varepsilon)\det T$. Substitution gives exactly the two displayed ratios. The image is measurable by Lemma \ref{lem:finite-volume-lipschitz-image-measurable}. Keeping the denominator factors explicit avoids reversing an inequality when comparing the Jacobian integral with $\det T\,|E|$.
\end{proof}

\begin{proposition}[The area identity on the full-rank differentiability set]\label{prop:area-full-rank-lipschitz}
For Lipschitz $f:\mathbb R^m\to\mathbb R^n$ and a finite-volume measurable set $A$ of full-rank differentiability points,
\[
 \int_AJ_mDf\,dx
 =\int_{\mathbb R^n}\#(A\cap f^{-1}(y))\,d\mathcal H^m(y).
\]
The multiplicity function is measurable and finite almost everywhere. This is the full-rank case of \cite{Sim18}, Theorem 3.3.
\end{proposition}
\begin{proof}
For a fixed $\varepsilon$, partition $A$ by Proposition \ref{prop:area-full-rank-metric-pieces}. Each restriction is injective, and each image is measurable by Lemma \ref{lem:finite-volume-lipschitz-image-measurable}. Therefore, at every target point,
\[
 \#(A\cap f^{-1}(y))=\sum_\ell\chi_{f(E_\ell)}(y).
\]
This proves measurability, including the value infinity when infinitely many pieces contribute. Sum Lemma \ref{lem:area-one-metric-piece} over the disjoint domain pieces. Monotone convergence identifies the image-measure sum with the multiplicity integral, so the same two comparison factors bound it by $\int_AJ_mDf$. This last integral is finite because $J_mDf\leq(\operatorname{Lip}f)^m$ almost everywhere and $|A|<\infty$. Let $\varepsilon\downarrow0$ to obtain equality. Finiteness of the integral gives finiteness of the multiplicity for almost every target point. The partitions may depend on $\varepsilon$; the function being integrated does not.
\end{proof}

\begin{lemma}[An injective graph regularization of a rank-deficient map]\label{lem:area-graph-regularization}
Let $f:\mathbb R^m\to\mathbb R^n$ be Lipschitz with constant $L$, and set $F_\varepsilon(x)=(f(x),\varepsilon x)$ for $0<\varepsilon\leq1$. This map is injective and Lipschitz. At every differentiability point of $f$, its derivative has rank $m$ and
\[
 J_mDF_\varepsilon
 =\sqrt{\det(Df^*Df+\varepsilon^2I_m)}.
\]
If $\operatorname{rank}Df<m$, this is at most $\varepsilon(L^2+1)^{(m-1)/2}$.
\end{lemma}
\begin{proof}
The last coordinate block determines $x$, proving injectivity, and $\operatorname{Lip}F_\varepsilon\leq\sqrt{L^2+\varepsilon^2}$. Its derivative has the vertical block $\varepsilon I_m$, so is injective. Its Gram matrix is $Df^*Df+\varepsilon^2I_m$. At a rank-deficient point at least one eigenvalue of $Df^*Df$ is zero, and all eigenvalues are at most $L^2$. Thus one regularized eigenvalue is $\varepsilon^2$, while the others are at most $L^2+1$. Taking the square root of their product gives the stated bound. Its order is $\varepsilon$ in the worst rank-$m-1$ case; a uniform quadratic error in $\varepsilon$ is unnecessary and generally false.
\end{proof}

\begin{proposition}[The critical differentiability set has a null image]\label{prop:area-critical-image-null}
Let $f:\mathbb R^m\to\mathbb R^n$ be Lipschitz and let $A$ be a measurable set of differentiability points with $\operatorname{rank}Df<m$. Then $\mathcal H^m(f(A))=0$.
\end{proposition}
\begin{proof}
First assume $|A|<\infty$. Lemma \ref{lem:area-graph-regularization} and Proposition \ref{prop:area-full-rank-lipschitz}, now in target $\mathbb R^{n+m}$, give
\[
 \mathcal H^m(F_\varepsilon(A))
 =\int_AJ_mDF_\varepsilon\,dx
 \leq\varepsilon(L^2+1)^{(m-1)/2}|A|.
\]
Here injectivity of $F_\varepsilon$ turns its multiplicity into the indicator of its image. Orthogonal projection onto the first $n$ coordinates is one-Lipschitz and maps $F_\varepsilon(A)$ onto $f(A)$. Consequently the same upper bound holds for $\mathcal H^m(f(A))$. Let $\varepsilon\downarrow0$. General $A$ is a countable union of finite-volume intersections with bounded balls, so its image is a countable union of null images. This proves the critical case using the already established full-rank case, rather than assuming the desired area formula on critical points.
\end{proof}

\begin{theorem}[Euclidean area with measurable multiplicity]\label{thm:euclidean-area-local-proof}
Let $V\subset\mathbb R^m$ be open, let $f:V\to\mathbb R^n$ be locally Lipschitz, and let $A\subset V$ be Lebesgue measurable. Then the counting function $N_A(y)=\#(A\cap f^{-1}(y))$ is $\mathcal H^m$-measurable, and
\[
 \int_AJ_mDf\,dx=\int_{\mathbb R^n}N_A(y)\,d\mathcal H^m(y).
\]
The Jacobian is assigned zero on the null set of nondifferentiability. Infinite values are permitted. This supplies the proof of \cite{Sim18}, Theorem 3.3, through its separate full-rank, critical and exceptional cases.
\end{theorem}
\begin{proof}
First let $f$ be globally Lipschitz and $|A|<\infty$. Partition $A$ into its full-rank differentiability part $R$, its rank-deficient differentiability part $C$, and its nondifferentiability part $Z$. Measurable derivative representatives make these measurable. Rademacher's theorem, Theorem \ref{thm:rademacher}, gives $|Z|=0$, so Lemma \ref{lem:finite-volume-lipschitz-image-measurable} gives $\mathcal H^m(f(Z))=0$. Proposition \ref{prop:area-critical-image-null} gives the same conclusion for $f(C)$. Their multiplicity functions vanish outside those null images and are measurable for complete Hausdorff measure, regardless of possibly infinite values on the images. Their Jacobian integrals are zero. Proposition \ref{prop:area-full-rank-lipschitz} treats $R$. Counts add on disjoint domain subsets, and this proves the identity on $A$.

For the local case, cover $V$ by countably many bounded balls with closures in larger balls on which $f$ is Lipschitz. Extend its coordinate functions from each larger ball to all of $\mathbb R^m$ by Proposition \ref{prop:mcshane-extension}. The vector extension is Lipschitz, possibly with a larger constant. On the smaller ball it agrees with $f$ on a neighborhood of every point, so its derivative agrees wherever that of $f$ exists. Partition $A$ disjointly among the smaller balls and apply the established finite-volume formula on each part. Sum by monotone convergence. This proves both the identity and the measurability of its counting function for arbitrary $A$, while preserving every local derivative. The proof uses neither Whitney extension nor rectifiable-set decomposition or coarea.
\end{proof}

\begin{corollary}[A null image is equivalent to a vanishing area Jacobian]\label{cor:area-null-image-criterion}
Under Theorem \ref{thm:euclidean-area-local-proof}, $\mathcal H^m(f(A))=0$ if and only if $J_mDf=0$ almost everywhere on $A$. In particular, a measurable positive-volume set on which the area Jacobian is positive almost everywhere has an image of positive $m$-dimensional measure.
\end{corollary}
\begin{proof}
The image is measurable by countable local versions of Lemma \ref{lem:finite-volume-lipschitz-image-measurable}. Its indicator is bounded above by $N_A$, since every point of the image has at least one preimage. If the Jacobian vanishes almost everywhere, the area identity gives $\int N_A=0$ and hence zero measure for the image. Conversely if the image is null, $N_A$ vanishes off a null set, so its integral and the Jacobian integral are zero. Nonnegativity of the Jacobian gives its almost-everywhere vanishing. The final assertion is the contrapositive with a positive integral, which may be infinite.
\end{proof}

\begin{proposition}[Finite-volume domains have finite multiplicity almost everywhere]\label{prop:area-integrable-multiplicity}
If $f:\mathbb R^m\to\mathbb R^n$ has Lipschitz constant $L$ and $A$ is measurable with finite volume, then
\[
 \int_{\mathbb R^n}N_A\,d\mathcal H^m\leq L^m|A|,
 \qquad
 \mathcal H^m(\{N_A\geq q\})\leq\frac{L^m|A|}{q}
 \quad(q=1,2,\ldots).
\]
In particular almost every target point has finitely many preimages in $A$. This conclusion permits exceptional infinite fibers.
\end{proposition}
\begin{proof}
All singular values of $Df$ are at most $L$, so $J_mDf\leq L^m$ at differentiability points. Theorem \ref{thm:euclidean-area-local-proof} gives the integral bound. On $\{N_A\geq q\}$ the counting function is at least $q$, giving the second inequality by integration. The set where $N_A=\infty$ is contained in every such superlevel set, so its measure is zero on letting $q\to\infty$. A constant map can still have an entire domain as the fiber at its one exceptional image point.
\end{proof}
